\documentclass{article}
\usepackage[utf8]{inputenc}
\usepackage[german]{babel}
\usepackage{amsmath, amssymb, amsthm}
\usepackage{booktabs}
\usepackage{array}

\newtheorem{definition}{Definition}
\newtheorem{theorem}{Satz}
\newcolumntype{L}[1]{>{\raggedright\arraybackslash}p{#1}}

\title{Vom Durchschnitt von Topologien zu $\sigma$-Algebren}
\author{Nach Daniel J. Velleman: \emph{How to Prove It}}
\date{\today}

\begin{document}

\maketitle

\section*{1. Definition einer Topologie}

Sei $X$ eine Grundmenge. Ein Mengensystem $\mathcal{T} \subseteq \mathcal{P}(X)$ heißt \textbf{Topologie} auf $X$, wenn die folgenden drei Axiome erfüllt sind:
\begin{enumerate}
    \item T1 (Leere Menge und Ganzraum): $\emptyset \in \mathcal{T} \land X \in \mathcal{T}$
    \item T2 (Endliche Durchschnitte): $\forall U, V \, \big( (U \in \mathcal{T} \land V \in \mathcal{T}) \rightarrow (U \cap V \in \mathcal{T}) \big)$
    \item T3 (Beliebige Vereinigungen): $\forall \mathcal{A} \, \big( \mathcal{A} \subseteq \mathcal{T} \rightarrow \bigcup \mathcal{A} \in \mathcal{T} \big)$
\end{enumerate}

---

\section*{2. Der Satz vom Durchschnitt einer Familie}

\begin{theorem}
Sei $X$ eine Menge und $I$ eine nichtleere Indexmenge. Für jedes $i \in I$ sei $\mathcal{T}_i$ eine Topologie auf $X$. Dann ist der Durchschnitt der Familie, definiert als
\[
\mathcal{T}_{\cap} = \bigcap_{i \in I} \mathcal{T}_i = \{ U \subseteq X \mid \forall i \in I \, (U \in \mathcal{T}_i) \},
\]
ebenfalls eine Topologie auf $X$.
\end{theorem}

---

\section*{3. Scratch Work (Exemplarisch für Axiom T2)}

Um zu beweisen, dass $\mathcal{T}_{\cap}$ eine Topologie ist, müssen wir alle drei Axiome zeigen. Schauen wir uns das Velleman-Scratch-Work für das oft kniffligste Axiom an: **T2 (Stabilität unter endlichen Durchschnitten)**.

\begin{center}
\small
\begin{tabular}{ L{0.45\textwidth} | L{0.45\textwidth} }
\toprule
\textbf{Givens} & \textbf{Goals} \\
\midrule
% Schritt 0: Ausgangssituation
1. $\forall i \in I \, (\mathcal{T}_i \text{ ist Topologie})$ \newline
2. $\mathcal{T}_{\cap} = \{ U \subseteq X \mid \forall i \in I \, (U \in \mathcal{T}_i) \}$ &
\emph{Goal:} $\forall U, V \big( (U \in \mathcal{T}_{\cap} \land V \in \mathcal{T}_{\cap}) \rightarrow U \cap V \in \mathcal{T}_{\cap} \big)$ \\[1ex]
\hline

% Schritt 1: Goal abbauen (Allquantoren und Implikation)
1. [Givens wie oben] \newline
3. Seien $U$ und $V$ beliebige Mengen. \newline
4. $U \in \mathcal{T}_{\cap}$ (Annahme) \newline
5. $V \in \mathcal{T}_{\cap}$ (Annahme) &
\emph{Neues Goal:} $U \cap V \in \mathcal{T}_{\cap}$ \newline
\newline
(\emph{Laut Definition von $\mathcal{T}_{\cap}$ bedeutet das:}) \newline
\emph{Goal*:} $\forall i \in I \, (U \cap V \in \mathcal{T}_i)$ \\[1ex]
\hline

% Schritt 2: Givens 4 und 5 mittels Definition von T_neu aufspalten
1. [Givens wie oben] \newline
3. $U, V$ beliebig \newline
6. $\forall i \in I \, (U \in \mathcal{T}_i)$ \quad (aus Given 4) \newline
7. $\forall i \in I \, (V \in \mathcal{T}_i)$ \quad (aus Given 5) &
\emph{Goal*:} $\forall i \in I \, (U \cap V \in \mathcal{T}_i)$ \\[1ex]
\hline

% Schritt 3: Allquantor im Goal angreifen
1. [Givens wie oben] \newline
3. $U, V$ beliebig; 6. und 7. wie oben \newline
8. Sei $i_0 \in I$ ein beliebiges, festes Indexelement. &
\emph{Neues Goal:} $U \cap V \in \mathcal{T}_{i_0}$ \\[1ex]
\hline

% Schritt 4: Givens mit dem konkreten i_0 instanziieren
3. $U, V$ beliebig; 8. $i_0 \in I$ beliebig \newline
9. $U \in \mathcal{T}_{i_0}$ \quad (Instanziierung von 6 mit $i_0$) \newline
10. $V \in \mathcal{T}_{i_0}$ \quad (Instanziierung von 7 mit $i_0$) \newline
11. $\mathcal{T}_{i_0}$ erfüllt T2 (da es eine Topologie ist!) &
\emph{Goal:} $U \cap V \in \mathcal{T}_{i_0}$ \\[1ex]
\hline

% Schritt 5: Axiom T2 von T_{i_0} nutzen (Modus Ponens)
% Wir instanziieren das Axiom T2 von T_{i_0} mit unseren Mengen U und V
3. $U,V$ beliebig; 8. $i_0$ beliebig; 9., 10. wie oben \newline
12. $(U \in \mathcal{T}_{i_0} \land V \in \mathcal{T}_{i_0}) \rightarrow U \cap V \in \mathcal{T}_{i_0}$ \newline
13. $U \cap V \in \mathcal{T}_{i_0}$ \quad (Modus Ponens aus 9, 10 und 12) &
\emph{Goal:} $U \cap V \in \mathcal{T}_{i_0}$ \newline
\newline
\checkmark \quad (Ziel erreicht!) \\
\bottomrule
\end{tabular}
\end{center}

\end{document}
